\title{The Era of 1-bit LLMs: All Large Language Models are in 1.58 Bits}

\begin{abstract}
Recent advancements, including BitNet~\cite{bitnet}, are ushering in a new generation of 1-bit Large Language Models (LLMs). In this study, we present a 1-bit LLM version called \textbf{\text{BitNet b1.58}}, where every parameter (or weight) of the LLM is ternary \{-1, 0, 1\}. This model achieves comparable performance to full-precision (i.e., FP16 or BF16) Transformer LLMs of equivalent model size and training data in terms of both perplexity and downstream task accuracy, while offering substantial improvements in latency, memory usage, throughput, and energy efficiency. More importantly, the 1.58-bit LLM introduces a novel scaling law and methodology for developing future LLMs that balance high performance with cost-effectiveness. Additionally, it facilitates a new computational paradigm and paves the way for designing dedicated hardware tailored to 1-bit LLMs.
\end{abstract}

\section{The Era of 1-bit LLMs}

In recent times, the AI field has experienced rapid expansion in both the scale and capabilities of Large Language Models (LLMs). These models have achieved impressive results across diverse natural language processing tasks. However, their growing size presents deployment challenges and raises environmental and financial concerns due to substantial energy demands.

A common strategy to mitigate these issues is post-training quantization, which produces low-bit models for inference~\cite{smoothquant, gptq, quip, quip_sharp}. This method lowers the precision of weights and activations, significantly decreasing the memory footprint and computational load of LLMs. The progression has been towards reducing precision from 16 bits down to lower bitwidths, such as 4-bit models~\cite{gptq, awq}. Nevertheless, despite its prevalence in industry-scale LLM deployments, post-training quantization remains a suboptimal approach.

Recent studies on 1-bit model architectures, such as BitNet~\cite{bitnet}, offer a promising approach to lowering the expenses associated with LLMs while preserving their effectiveness. Standard LLMs operate with 16-bit floating-point values (e.g., FP16 or BF16), and their primary computation involves matrix multiplication. Consequently, the main computational burden arises from floating-point addition and multiplication operations. In contrast, BitNet’s matrix multiplication relies solely on integer addition, which greatly reduces the energy consumption required by LLMs. Since power consumption is often the limiting factor for compute performance in many chips, these energy savings can also lead to accelerated computation.

Beyond computation, transferring model parameters from DRAM to the on-chip accelerator's memory (e.g., SRAM) can be costly during inference. Attempts to increase SRAM size to boost throughput have been made, but these come with substantially higher costs compared to DRAM. Relative to full-precision models, 1-bit LLMs possess a significantly smaller memory footprint in terms of both capacity and bandwidth. This reduction can greatly cut down the cost and duration needed to load weights from DRAM, resulting in faster and more efficient inference.

In this paper, we propose a notable 1-bit LLM variant named \textbf{\text{BitNet b1.58}}, where every parameter is ternary and takes values from \{-1, 0, 1\}. We introduced the additional zero value to the original 1-bit BitNet, yielding 1.58 bits in binary representation. \text{BitNet b1.58} maintains all advantages of the original 1-bit BitNet, including its novel computation scheme that almost eliminates multiplication operations in matrix multiplication and allows for high optimization. It also matches the original 1-bit BitNet in energy consumption and surpasses FP16 LLM baselines in memory efficiency, throughput, and latency. Moreover, \text{BitNet b1.58} provides two extra benefits. First, its modeling power is enhanced thanks to explicit feature filtering enabled by the zero value in the model weights, which significantly boosts 1-bit LLM performance. Second, our experiments indicate that \text{BitNet b1.58}, starting from a 3B model size and under identical configurations (e.g., model size, training tokens), can achieve performance on par with full precision (i.e., FP16) baselines in both perplexity and downstream tasks.

\section{BitNet b1.58}

\text{BitNet b1.58} builds upon the BitNet framework, which is a Transformer model that substitutes \emph{nn.Linear} layers with \emph{BitLinear}. It is trained from the ground up using 1.58-bit precision for weights and 8-bit precision for activations. Compared to the original BitNet, several modifications have been introduced, which we outline below.

\paragraph{Quantization Function.} To restrict the weights to values in \{-1, 0, +1\}, we utilize an \emph{absmean} quantization method. This approach first normalizes the weight matrix by its mean absolute value, then rounds each element to the nearest integer within \{-1, 0, +1\}:

\begin{equation}
    \widetilde{W} = \mathrm{RoundClip}(\frac{W}{\gamma + \epsilon}, -1, 1),
\end{equation}

\begin{equation}
    \mathrm{RoundClip}(x, a, b) = \max(a, \min(b, \mathrm{round}(x))),
\end{equation}

\begin{equation}
    \gamma = \frac{1}{nm}\sum_{ij} |W_{ij}|.
\end{equation}

For activation quantization, we maintain the same method as BitNet, except that we no longer scale activations before the nonlinearities to the range $[0, Q_b]$. Instead, activations are scaled per token to lie within $[-Q_b, Q_b]$, removing the need for zero-point quantization. This adjustment simplifies implementation and system-level optimization, while having a negligible impact on performance based on our experimental observations.

\paragraph{LLaMA-alike Components.}

The LLaMA~\cite{llama, llama2} architecture has become the standard backbone for many open-source large language models. To align with the open-source ecosystem, \text{BitNet b1.58} incorporates LLaMA-like components. Specifically, it employs RMSNorm~\cite{rmsnorm}, SwiGLU~\cite{swiglu}, rotary positional embeddings~\cite{rope}, and eliminates all bias terms. This design choice enables \text{BitNet b1.58} to be easily integrated into widely-used open-source platforms such as Huggingface, vLLM~\cite{vllm}, and llama.cpp\footnote{\url{https://github.com/ggerganov/llama.cpp}} with minimal modification efforts.

\section{Results}

We conducted a comparison between \text{BitNet b1.58} and our reproduced FP16 \text{LLaMA LLM} across various model sizes. To guarantee an equitable comparison, both models were pre-trained on the RedPajama dataset~\cite{redpajama} with 100 billion tokens. Their zero-shot capabilities were assessed on multiple language tasks, including ARC-Easy~\cite{arc}, ARC-Challenge~\cite{arc}, Hellaswag~\cite{hellaswag}, Winogrande~\cite{winoGrande}, PIQA~\cite{piqa}, OpenbookQA~\cite{openbookqa}, and BoolQ~\cite{boolq}. Additionally, we reported validation perplexity on the WikiText2~\cite{wikitext2} and C4~\cite{c4} datasets.

We evaluated the runtime GPU memory consumption and latency for both \text{LLaMA LLM} and \text{BitNet b1.58}. These measurements were performed using the FasterTransformer\footnote{\url{https://github.com/NVIDIA/FasterTransformer}} codebase, which is optimized for minimizing inference latency of LLMs on GPU platforms. The 2-bit kernel developed by Ladder~\cite{ladder} was also integrated for \text{BitNet b1.58}. We reported the per-output-token processing time, as it constitutes the primary inference cost.

Table~\ref{tab:ppl} presents a summary of perplexity and computational cost for both \text{BitNet b1.58} and \text{LLaMA LLM}. It indicates that \text{BitNet b1.58} begins to match the full precision \text{LLaMA LLM} perplexity at a 3B model size, while delivering 2.71 times faster speed and requiring 3.55 times less GPU memory. Notably, \text{BitNet b1.58} with a 3.9B parameter model is 2.4 times faster, uses 3.32 times less memory, yet significantly outperforms the \text{LLaMA LLM} 3B model.

Detailed zero-shot accuracy outcomes on the downstream tasks are shown in Table~\ref{tab:zero-shot}. We adopted the evaluation pipeline from \emph{lm-evaluation-harness}\footnote{\url{https://github.com/EleutherAI/lm-evaluation-harness}}. The findings demonstrate that the performance disparity between \text{BitNet b1.58} and \text{LLaMA LLM} diminishes with increasing model size. More importantly, \text{BitNet b1.58} achieves parity with the full precision baseline starting at the 3B scale. Consistent with the perplexity results, the downstream task evaluations confirm that \text{BitNet b1.58} 3.9B surpasses \text{LLaMA LLM} 3B while incurring lower memory usage and latency. This establishes \text{BitNet b1.58} as a Pareto improvement over current state-of-the-art LLMs.

\paragraph{Memory and Latency}

We additionally increased the model size to 7B, 13B, and 70B and assessed the associated cost. Figure~\ref{fig:latency-memory-energy} displays the patterns of latency and memory usage, demonstrating that the speed-up improves as the model size increases. Notably, \text{BitNet b1.58} 70B operates 4.1 times faster than the \text{LLaMA LLM} baseline. This improvement arises because the time required for \emph{nn.Linear} grows with the model size. Memory usage exhibits a similar pattern, since the embedding stays at full precision and its share of memory decreases for larger models. Both latency and memory were evaluated using a 2-bit kernel, indicating there is potential for further optimization to reduce costs.

\paragraph{Energy}

We also estimated the energy consumption for arithmetic operations of both \text{BitNet b1.58} and \text{LLaMA LLM}. Our main focus was on matrix multiplication calculations, as these dominate the computational cost in LLMs. Figure~\ref{fig:energy} presents the breakdown of energy costs. Most of the \text{BitNet b1.58} consumption arises from INT8 addition operations, whereas \text{LLaMA LLM} includes both FP16 addition and FP16 multiplication. Based on the energy model from~\cite{energycost, pokebnn}, \text{BitNet b1.58} reduces energy consumption for matrix multiplication arithmetic operations by 71.4 times on 7nm chips. We also reported the total end-to-end energy cost for models processing 512 tokens. Our findings indicate that as model size grows, \text{BitNet b1.58} becomes progressively more energy-efficient compared to the FP16 \text{LLaMA LLM} baseline. This is because the proportion of \emph{nn.Linear} increases with model size, while the relative cost from other components decreases for larger models.

\paragraph{Throughput}

We evaluate the throughput of \text{BitNet b1.58} compared to the \text{LLaMA LLM} with 70B parameters using two 80GB A100 GPUs, employing pipeline parallelism~\cite{gpipe} to enable running the \text{LLaMA LLM} 70B on these devices. We increased the batch size progressively until the GPU memory capacity was exhausted, with a fixed sequence length of 512. As shown in Table~\ref{tab:throughput}, \text{BitNet b1.58} 70B supports batch sizes up to 11 times larger than \text{LLaMA LLM}, delivering an 8.9-fold increase in throughput.

\textbf{\text{BitNet b1.58} establishes a new scaling relationship between model performance and inference cost}. For reference, based on the outcomes in Figures~\ref{fig:latency-memory-energy} and \ref{fig:energy}, we observe the following equivalences between models of different sizes in 1.58-bit and 16-bit precision:

\begin{itemize}
    \item A 13B \text{BitNet b1.58} model is more efficient regarding latency, memory usage, and energy consumption than a 3B FP16 LLM.
    \item A 30B \text{BitNet b1.58} model is more efficient in terms of latency, memory, and energy compared to a 7B FP16 LLM.
    \item A 70B \text{BitNet b1.58} model is more efficient considering latency, memory, and energy than a 13B FP16 LLM.
\end{itemize}

\paragraph{Training with 2T Tokens}

The quantity of training tokens plays a vital role for LLMs. To evaluate the scalability of \text{BitNet b1.58} concerning token count, we trained a \text{BitNet b1.58} model on 2T tokens using the data pipeline from StableLM-3B~\cite{StableLM-3B-4E1T}, which represents the current state-of-the-art open-source 3B model. Both models were assessed on a benchmark comprising Winogrande~\cite{winoGrande}, PIQA~\cite{piqa}, SciQ~\cite{sciq}, LAMBADA~\cite{lambada}, and ARC-easy~\cite{arc}. We present the zero-shot accuracy results in Table~\ref{tab:2t}. For tasks measured with both accuracy and normalized accuracy, we computed the average of the two metrics. The performance data for StableLM 3B at 2T tokens was taken directly from its official technical report. Our results indicate that \text{BitNet b1.58} attains better performance across all evaluated tasks, demonstrating that 1.58-bit LLMs possess strong generalization abilities as well.

\section{Discussion and Future Work}

\textbf{1-bit Mixture-of-Experts (MoE) LLMs}

Mixture-of-Experts (MoE) has emerged as a cost-efficient strategy for LLMs. Although it drastically lowers computation FLOPs, the large memory requirements and inter-chip communication bottlenecks pose challenges for its implementation and usage. These issues can be mitigated through 1.58-bit LLMs. Primarily, the smaller memory footprint decreases the number of devices needed to host MoE models. Additionally, it greatly diminishes the cost associated with transferring activations over networks. Ultimately, if the entire model can reside on a single chip, these overheads would be eliminated.

\textbf{Native Support of Long Sequence in LLMs}

With the rise of LLMs, managing long sequences has become increasingly important. A key obstacle in long sequence inference is the memory usage caused by KV caches. \text{BitNet b1.58} marks a crucial advance toward native long sequence support by halving the activation bit-width from 16 to 8 bits, effectively enabling the context length to double using the same hardware resources. Furthermore, this can be losslessly compressed to 4 bits or even less for 1.58-bit LLMs, which we identify as a direction for future exploration.

\textbf{LLMs on Edge and Mobile}

Adopting 1.58-bit LLMs holds substantial promise for enhancing language model performance on edge and mobile platforms. These devices are frequently constrained by limited memory and computational capabilities, which restrict the complexity and effectiveness of LLMs. By lowering memory and energy demands, 1.58-bit LLMs facilitate deployment on such devices, opening the door to numerous applications previously unattainable. This advancement can significantly boost the functionality of edge and mobile hardware and foster novel applications of LLMs. Moreover, 1.58-bit LLMs are better suited to CPU architectures, the predominant processors in these devices, allowing \text{BitNet b1.58} to run efficiently and further elevate their performance and utility.

\textbf{New Hardware for 1-bit LLMs}

Recent initiatives like~Groq\footnote{\url{https://groq.com/}} have showcased encouraging outcomes and considerable promise for developing specialized hardware (e.g., LPUs) tailored to LLMs. Building on this progress, we envision and urge efforts toward creating new hardware and system architectures specifically optimized for 1-bit LLMs, leveraging the novel computational framework introduced by BitNet~\cite{bitnet}.

\end{document}